\documentclass[12pt,a4paper]{article}

\usepackage[utf8]{inputenc}
\usepackage[T1]{fontenc}
\usepackage[brazil]{babel}
\usepackage{lmodern}
\usepackage{geometry}
\usepackage{microtype}
\usepackage{array}
\usepackage{tabularx}
\usepackage{enumitem}
\usepackage{xcolor}
\usepackage{fancyhdr}
\usepackage{amsmath}

\geometry{top=1.7cm,bottom=1.7cm,left=1.8cm,right=1.8cm}
\setlength{\parindent}{0pt}
\setlength{\parskip}{0.45em}
\setlist[itemize]{leftmargin=*,itemsep=0.25em,topsep=0.2em}
\setlist[enumerate]{leftmargin=*,itemsep=0.35em,topsep=0.2em}
\definecolor{azul}{HTML}{173B57}
\definecolor{cinza}{HTML}{F2F4F7}
\definecolor{cinzaescuro}{HTML}{5E6670}

\pagestyle{fancy}
\setlength{\headheight}{15pt}
\fancyhf{}
\lhead{Sequência didática de \textit{machine learning}}
\rhead{Aula 2: regressão linear}
\cfoot{\thepage}
\renewcommand{\headrulewidth}{0.3pt}

\newcommand{\campo}[1]{\textbf{#1}}
\newcommand{\linha}{\rule{\linewidth}{0.35pt}}
\newcommand{\bloco}[1]{%
  \vspace{0.25em}
  \noindent\colorbox{cinza}{\parbox{0.96\linewidth}{#1}}%
  \vspace{0.3em}
}

\begin{document}

\begin{center}
  {\LARGE\bfseries\color{azul} Aula 2: regressão linear}\par
  \vspace{0.25em}
  {\large Massa, alongamento da mola, previsão e erro}\par
  \vspace{0.35em}
  {\small Roteiro de aplicação da sequência didática}\par
\end{center}

\bloco{\textbf{Lugar desta aula na pesquisa.} Este é o segundo encontro da sequência de oito aulas. A atividade introduz a regressão linear como forma de estimar um valor numérico a partir de exemplos. O registro individual ao final integra o conjunto de evidências previsto para os encontros 2, 4 e 6.}

\section*{1. Identificação}

\begin{tabularx}{\linewidth}{@{}>{\bfseries}p{4.1cm}X@{}}
Turma: & \rule{8.5cm}{0.35pt} \\
Data: & \rule{8.5cm}{0.35pt} \\
Duração: & 50 minutos \\
Organização: & Exploração individual; registro individual ao final. \\
Recursos: & Plataforma no módulo de regressão linear, computadores do laboratório ou Chromebooks, quadro e folha de registro. \\
\end{tabularx}

\section*{2. Objetivos da aula}

Ao final do encontro, espera-se que os estudantes consigam reconhecer uma situação de previsão de valor numérico, interpretar uma tendência em gráfico de dispersão e relacionar uma alteração na inclinação da reta ou nos dados a mudanças na previsão e no erro.

\textbf{Relação curricular:} EM13CO11 e EM13CO12.

\textbf{Conceitos em foco:} dado, variável, previsão numérica, gráfico de dispersão, reta de tendência, inclinação, valor observado, valor previsto, resíduo e erro.

\section*{3. Preparação antes da aula}

\begin{itemize}
  \item Abrir o módulo de regressão linear e confirmar que o cenário de massa e alongamento da mola está disponível.
  \item Conferir se os eixos, os pontos, a reta e os controles podem ser vistos pela turma.
  \item Separar uma folha de registro individual para cada estudante e usar o mesmo código anônimo adotado no diagnóstico inicial.
  \item Confirmar que cada estudante terá acesso individual ao módulo, em computador do laboratório ou Chromebook.
\end{itemize}

\bloco{\textbf{Cuidado didático.} A situação da mola é um cenário sintético. A finalidade é observar relações entre valores, dados e modelo, sem tratar o gráfico como previsão exata de um fenômeno real.}

\section*{4. Desenvolvimento da aula}

\renewcommand{\arraystretch}{1.25}
\begin{tabularx}{\linewidth}{|>{\centering\arraybackslash}p{1.65cm}|p{5.0cm}|X|}
\hline
\textbf{Tempo} & \textbf{Ação com a turma} & \textbf{Mediação docente} \\
\hline
0--7 min & Retomar a aula anterior: dados entram em um sistema e uma saída é produzida. Apresentar a diferença entre escolher uma categoria e estimar um valor. & Perguntar: ``Se queremos saber o alongamento para uma massa ainda não testada, qual valor queremos prever?'' Registrar no quadro que a saída será um número. \\
\hline
7--12 min & Apresentar a tabela e o gráfico de massa aplicada e alongamento da mola. Identificar os eixos, os pontos observados e o valor a estimar. & Nomear cada eixo antes de apresentar a reta. Perguntar qual variável representa a massa e qual representa o alongamento. \\
\hline
12--18 min & Pedir que cada estudante antecipe o que acontece com a previsão de um ponto à direita quando a inclinação da reta aumenta. & Solicitar uma resposta específica: aumenta, diminui ou permanece igual? Pedir que indiquem no desenho por que pensam isso. \\
\hline
18--33 min & Explorar o módulo de forma guiada. Alterar a inclinação, observar a reta e comparar valores previstos e observados. Se houver tempo, mostrar um passo de ajuste da reta. & Apresentar o resíduo como diferença entre o ponto observado e o valor indicado pela reta. Não exigir cálculo geral de erro nem ajuste manual perfeito. \\
\hline
33--43 min & Retomar previsões registradas individualmente e confrontá-las com o gráfico. Explorar um ponto que fique distante da reta. & Perguntar o que muda na previsão quando a reta se move e o que um ponto distante informa sobre a diferença entre observado e previsto. \\
\hline
43--50 min & Aplicar e recolher o registro individual. Encerrar com uma síntese oral. & Esclarecer palavras do enunciado sem sugerir respostas. Reforçar que o registro não vale nota e interessa como explicação do estudante. \\
\hline
\end{tabularx}

\section*{5. Perguntas para conduzir a conversa}

\begin{itemize}
  \item ``O que os pontos do gráfico representam?''
  \item ``Qual valor a reta tenta estimar?''
  \item ``O que tende a acontecer com um ponto à direita quando a inclinação aumenta?''
  \item ``Onde aparece a diferença entre o valor observado e o valor previsto?''
  \item ``Uma reta próxima de alguns pontos garante que ela explique todos os casos?''
\end{itemize}

\section*{6. Evidência a observar}

O registro individual permite verificar se o estudante relaciona a inclinação da reta a uma mudança na previsão e se interpreta o resíduo como diferença entre valor observado e valor previsto. Também convém registrar, após a aula, dificuldades recorrentes da turma, ajustes de linguagem realizados e condições de infraestrutura que tenham interferido na atividade.

\newpage

\begin{center}
  {\LARGE\bfseries\color{azul} Registro individual}\par
  \vspace{0.2em}
  {\large Aula 2: massa, alongamento e previsão}\par
  \vspace{0.35em}
  {\small Material para recolhimento ao final da aula}\par
\end{center}

\bloco{\textbf{Orientação ao estudante.} Esta atividade não vale nota. Responda individualmente e explique como você pensou, mesmo que não tenha certeza. Use o gráfico e o que foi discutido na aula.}

\textbf{Código do participante:} \rule{5cm}{0.4pt}
\hfill
\textbf{Data:} \rule{3cm}{0.4pt}

\section*{Situação}

Em um experimento fictício, foram colocadas massas diferentes em uma mola. Para cada massa, foi observado quanto a mola aumentou de tamanho. Os pontos do gráfico mostram as observações, e a reta indica uma previsão de alongamento para cada massa.

\section*{Questão 1}

Para uma massa maior, localizada à direita do gráfico, a previsão de alongamento tende a aumentar ou diminuir quando a inclinação da reta aumenta? Explique usando a posição da reta no gráfico.

\vspace{0.25em}
\linha\par\vspace{0.75em}
\linha

\section*{Questão 2}

Um ponto observado ficou acima da reta. O que isso indica sobre o alongamento observado e o alongamento previsto pela reta para aquela massa?

\vspace{0.25em}
\linha\par\vspace{0.75em}
\linha\par\vspace{0.75em}
\linha

\section*{Questão 3}

Uma nova massa foi colocada na mola e a reta indicou um alongamento de 12 cm. Depois do experimento, foi observado um alongamento de 15 cm. Qual é a diferença entre os dois valores? O que essa diferença mostra sobre a previsão?

\vspace{0.25em}
\linha\par\vspace{0.75em}
\linha\par\vspace{0.75em}
\linha

\end{document}
